% =================================================
% 文件: Kubernetes.tex
% 章节: Kubernetes 容器编排
% 版本: v1.0
% =================================================

\documentclass[12pt,UTF8]{ctexbook}

\input{../common/page}

\xeCJKsetup{AutoFallBack=true}
\setCJKfamilyfont{hei}{SimHei}
\setCJKfamilyfont{kai}{KaiTi}

\usepackage{titletoc}
\titlecontents{chapter}[0pt]{\vspace{3mm}\bf\addvspace{2pt}\filright}{\contentspush{\thecontentslabel\hspace{0.8em}}}{}{\titlerule*[8pt]{.}\contentspage}

\usepackage[colorlinks,linkcolor=blue,citecolor=blue,urlcolor=blue]{hyperref}

\ctexset{
	part/name= {第,卷},
	part/number={\chinese{part}},
	chapter/name={第,章},
	chapter/number={\arabic{chapter}}
}

\usepackage{graphicx}
\graphicspath{{Images/}}

\usepackage[perpage]{footmisc}

\usepackage{enumitem}

\title{\heiti\zihao{0} Kubernetes 容器编排入门指南}
\author{WangFei}
\date{\today}

\begin{document}

\maketitle
\tableofcontents

\frontmatter
\chapter{前言}

本指南面向初学者，详细介绍 Kubernetes 的基础概念、核心组件、部署方法、资源管理、服务发现等知识。

\mainmatter

\chapter{Kubernetes 基础概念}
\label{chap:k8s_basics}

\section{什么是 Kubernetes}
\label{sec:what_is_k8s}

Kubernetes（简称K8s）是一个开源的容器编排平台，用于自动化部署、扩展和管理容器化应用。

Kubernetes 的主要功能：
\begin{itemize}
    \item \textbf{自动化部署}：自动部署应用容器
    \item \textbf{自动扩展}：根据负载自动增减容器数量
    \item \textbf{自愈能力}：自动重启失败的容器
    \item \textbf{服务发现}：自动发现和路由服务
    \item \textbf{负载均衡}：在多个容器间分配流量
\end{itemize}

\section{容器化基础}
\label{sec:containerization}

容器是一种轻量级的虚拟化技术，将应用及其依赖打包在一起。

容器与虚拟机的区别：
\begin{table}[htbp]
    \centering
    \caption{容器与虚拟机对比}
    \begin{tabular}{|c|p{5cm}|p{5cm}|}
        \hline
        \textbf{特性} & \textbf{容器} & \textbf{虚拟机} \\
        \hline
        隔离级别 & 进程级隔离 & 系统级隔离 \\
        \hline
        资源占用 & 低 & 高 \\
        \hline
        启动速度 & 秒级 & 分钟级 \\
        \hline
        镜像大小 & 小 & 大 \\
        \hline
        兼容性 & 需要兼容内核 & 完全独立 \\
        \hline
    \end{tabular}
\end{table}

\section{Kubernetes 架构}
\label{sec:k8s_architecture}

Kubernetes 集群由控制平面和工作节点组成：

\begin{table}[htbp]
    \centering
    \caption{Kubernetes 组件}
    \begin{tabular}{|c|p{8cm}|}
        \hline
        \textbf{组件} & \textbf{说明} \\
        \hline
        API Server & 集群的控制入口，处理所有请求 \\
        \hline
        etcd & 分布式键值存储，保存集群状态 \\
        \hline
        Scheduler & 调度器，将 Pod 分配到节点 \\
        \hline
        Controller Manager & 控制器管理器，维护期望状态 \\
        \hline
        Kubelet & 节点代理，管理 Pod 生命周期 \\
        \hline
        Kube-proxy & 网络代理，维护网络规则 \\
        \hline
    \end{tabular}
\end{table}

\chapter{核心概念}
\label{chap:core_concepts}

\section{Pod}
\label{sec:pod}

Pod 是 Kubernetes 的最小部署单元，包含一个或多个容器。

Pod 的特点：
\begin{itemize}
    \item 多个容器共享网络和存储
    \item 容器间通过 localhost 通信
    \item Pod 是原子性的，一起创建和销毁
\end{itemize}

Pod 的类型：
\begin{itemize}
    \item \textbf{自主 Pod}：独立运行的 Pod
    \item \textbf{控制器管理的 Pod}：由控制器管理的 Pod
\end{itemize}

\section{ReplicaSet}
\label{sec:replicaset}

ReplicaSet 确保指定数量的 Pod 副本在运行。

ReplicaSet 的功能：
\begin{itemize}
    \item 维持指定数量的 Pod 副本
    \item 自动替换失败的 Pod
    \item 支持滚动更新
\end{itemize}

\section{Deployment}
\label{sec:deployment}

Deployment 是更高层的抽象，管理 ReplicaSet 和 Pod。

Deployment 的功能：
\begin{itemize}
    \item 声明式管理 Pod
    \item 支持滚动更新和回滚
    \item 支持暂停和继续部署
\end{itemize}

\section{Service}
\label{sec:service}

Service 为一组 Pod 提供稳定的访问入口。

Service 的类型：
\begin{table}[htbp]
    \centering
    \caption{Service 类型}
    \begin{tabular}{|c|p{8cm}|}
        \hline
        \textbf{类型} & \textbf{说明} \\
        \hline
        ClusterIP & 集群内部 IP，默认类型 \\
        \hline
        NodePort & 在节点端口上暴露服务 \\
        \hline
        LoadBalancer & 通过负载均衡器暴露服务 \\
        \hline
        ExternalName & 映射到外部服务 \\
        \hline
    \end{tabular}
\end{table}

\section{Namespace}
\label{sec:namespace}

Namespace 用于隔离集群资源。

Namespace 的用途：
\begin{itemize}
    \item 隔离不同团队或项目的资源
    \item 控制资源访问权限
    \item 资源配额管理
\end{itemize}

\chapter{安装与配置}
\label{chap:installation}

\section{安装方式}
\label{sec:installation_methods}

Kubernetes 的安装方式：
\begin{table}[htbp]
    \centering
    \caption{安装方式}
    \begin{tabular}{|c|p{8cm}|}
        \hline
        \textbf{方式} & \textbf{说明} \\
        \hline
        kubeadm & 官方工具，快速部署 \\
        \hline
        minikube & 本地开发环境 \\
        \hline
        kind & 使用 Docker 运行集群 \\
        \hline
        云厂商托管 & AWS EKS、Azure AKS、GCP GKE \\
        \hline
    \end{tabular}
\end{table}

\section{kubeadm 安装}
\label{sec:kubeadm_install}

使用 kubeadm 安装 Kubernetes（CentOS Stream 9）：

\textbf{步骤1：准备环境}

禁用防火墙和 SELinux：
\begin{verbatim}
systemctl stop firewalld && systemctl disable firewalld
setenforce 0
sed -i 's/^SELINUX=.*/SELINUX=disabled/' /etc/selinux/config
\end{verbatim}

关闭 swap：
\begin{verbatim}
swapoff -a
sed -i '/swap/d' /etc/fstab
\end{verbatim}

设置内核参数：
\begin{verbatim}
cat > /etc/sysctl.d/k8s.conf <<EOF
net.bridge.bridge-nf-call-ip6tables = 1
net.bridge.bridge-nf-call-iptables = 1
EOF
sysctl --system
\end{verbatim}

\textbf{步骤2：安装 Docker}

\begin{verbatim}
dnf install -y dnf-utils
dnf config-manager --add-repo https://download.docker.com/linux/centos/docker-ce.repo
dnf install -y docker-ce docker-ce-cli containerd.io
systemctl enable --now docker
\end{verbatim}

\textbf{步骤3：安装 kubeadm、kubelet、kubectl}

\begin{verbatim}
cat > /etc/yum.repos.d/kubernetes.repo <<EOF
[kubernetes]
name=Kubernetes
baseurl=https://packages.cloud.google.com/yum/repos/kubernetes-el7-x86_64
enabled=1
gpgcheck=1
repo_gpgcheck=1
gpgkey=https://packages.cloud.google.com/yum/doc/yum-key.gpg https://packages.cloud.google.com/yum/doc/rpm-package-key.gpg
EOF

dnf install -y kubelet kubeadm kubectl
systemctl enable --now kubelet
\end{verbatim}

\textbf{步骤4：初始化集群}

\begin{verbatim}
kubeadm init --pod-network-cidr=10.244.0.0/16
\end{verbatim}

\textbf{步骤5：配置 kubectl}

\begin{verbatim}
mkdir -p $HOME/.kube
cp -i /etc/kubernetes/admin.conf $HOME/.kube/config
chown $(id -u):$(id -g) $HOME/.kube/config
\end{verbatim}

\textbf{步骤6：安装网络插件}

\begin{verbatim}
kubectl apply -f https://raw.githubusercontent.com/flannel-io/flannel/master/Documentation/kube-flannel.yml
\end{verbatim}

\chapter{资源管理}
\label{chap:resource_management}

\section{资源请求与限制}
\label{sec:resource_requests}

Pod 可以指定资源请求和限制：

\begin{verbatim}
apiVersion: v1
kind: Pod
metadata:
  name: resource-pod
spec:
  containers:
  - name: app
    image: nginx
    resources:
      requests:
        memory: "256Mi"
        cpu: "250m"
      limits:
        memory: "512Mi"
        cpu: "500m"
\end{verbatim}

\section{资源配额}
\label{sec:resource_quota}

Namespace 级别的资源配额：

\begin{verbatim}
apiVersion: v1
kind: ResourceQuota
metadata:
  name: ns-quota
spec:
  hard:
    requests.cpu: "4"
    requests.memory: 8Gi
    limits.cpu: "8"
    limits.memory: 16Gi
    pods: "20"
\end{verbatim}

\section{LimitRange}
\label{sec:limitrange}

限制 Pod 的默认资源使用：

\begin{verbatim}
apiVersion: v1
kind: LimitRange
metadata:
  name: default-limit
spec:
  limits:
  - default:
      cpu: "500m"
      memory: "512Mi"
    defaultRequest:
      cpu: "250m"
      memory: "256Mi"
    type: Container
\end{verbatim}

\chapter{服务发现与网络}
\label{chap:service_discovery}

\section{DNS 服务}
\label{sec:dns_service}

Kubernetes 内置 DNS 服务，Pod 可以通过域名访问服务。

DNS 解析格式：
\begin{itemize}
    \item Service：`<service-name>.<namespace>.svc.cluster.local`
    \item Pod：`<pod-ip>.pod.cluster.local`
\end{itemize}

\section{Ingress}
\label{sec:ingress}

Ingress 管理外部访问集群服务的规则。

Ingress 配置示例：
\begin{verbatim}
apiVersion: networking.k8s.io/v1
kind: Ingress
metadata:
  name: example-ingress
spec:
  rules:
  - host: example.com
    http:
      paths:
      - path: /
        pathType: Prefix
        backend:
          service:
            name: web-service
            port:
              number: 80
\end{verbatim}

\section{网络策略}
\label{sec:network_policy}

Network Policy 控制 Pod 间的网络流量。

\begin{verbatim}
apiVersion: networking.k8s.io/v1
kind: NetworkPolicy
metadata:
  name: deny-all
spec:
  podSelector: {}
  policyTypes:
  - Ingress
  - Egress
\end{verbatim}

\chapter{存储管理}
\label{chap:storage}

\section{存储类型}
\label{sec:storage_types}

Kubernetes 支持多种存储类型：

\begin{table}[htbp]
    \centering
    \caption{存储类型}
    \begin{tabular}{|c|p{8cm}|}
        \hline
        \textbf{类型} & \textbf{说明} \\
        \hline
        EmptyDir & 临时存储，Pod 删除时数据丢失 \\
        \hline
        HostPath & 使用节点本地文件系统 \\
        \hline
        PersistentVolume & 持久化存储卷 \\
        \hline
        PersistentVolumeClaim & 存储卷申请 \\
        \hline
        CSI & 容器存储接口 \\
        \hline
    \end{tabular}
\end{table}

\section{PersistentVolume}
\label{sec:pv}

创建 PersistentVolume：
\begin{verbatim}
apiVersion: v1
kind: PersistentVolume
metadata:
  name: pv-example
spec:
  capacity:
    storage: 10Gi
  accessModes:
    - ReadWriteOnce
  hostPath:
    path: /data/pv
\end{verbatim}

\section{PersistentVolumeClaim}
\label{sec:pvc}

创建 PersistentVolumeClaim：
\begin{verbatim}
apiVersion: v1
kind: PersistentVolumeClaim
metadata:
  name: pvc-example
spec:
  accessModes:
    - ReadWriteOnce
  resources:
    requests:
      storage: 5Gi
\end{verbatim}

\chapter{安全}
\label{chap:security}

\section{RBAC}
\label{sec:rbac}

基于角色的访问控制（RBAC）：

\begin{verbatim}
apiVersion: rbac.authorization.k8s.io/v1
kind: Role
metadata:
  namespace: default
  name: pod-reader
rules:
- apiGroups: [""]
  resources: ["pods"]
  verbs: ["get", "watch", "list"]
\end{verbatim}

创建 RoleBinding：
\begin{verbatim}
apiVersion: rbac.authorization.k8s.io/v1
kind: RoleBinding
metadata:
  name: read-pods
subjects:
- kind: User
  name: jane
roleRef:
  kind: Role
  name: pod-reader
  apiGroup: rbac.authorization.k8s.io
\end{verbatim}

\section{Secret}
\label{sec:secret}

Secret 用于存储敏感信息：

\begin{verbatim}
apiVersion: v1
kind: Secret
metadata:
  name: db-secret
type: Opaque
data:
  username: YWRtaW4=
  password: cGFzc3dvcmQ=
\end{verbatim}

\section{ServiceAccount}
\label{sec:serviceaccount}

ServiceAccount 为 Pod 提供身份：

\begin{verbatim}
apiVersion: v1
kind: ServiceAccount
metadata:
  name: app-sa
\end{verbatim}

\chapter{监控与日志}
\label{chap:monitoring}

\section{监控指标}
\label{sec:metrics}

使用 metrics-server 收集监控指标：

\begin{verbatim}
kubectl apply -f https://github.com/kubernetes-sigs/metrics-server/releases/latest/download/components.yaml
\end{verbatim}

查看节点资源使用：
\begin{verbatim}
kubectl top nodes
kubectl top pods
\end{verbatim}

\section{日志管理}
\label{sec:logging}

查看 Pod 日志：
\begin{verbatim}
kubectl logs <pod-name>
kubectl logs <pod-name> -f
kubectl logs <pod-name> --previous
\end{verbatim}

\section{告警}
\label{sec:alerting}

使用 Prometheus 和 Alertmanager 进行告警：

\begin{verbatim}
helm install prometheus prometheus-community/prometheus
helm install alertmanager prometheus-community/alertmanager
\end{verbatim}

\chapter{实践操作}
\label{chap:practice}

\section{部署应用}
\label{sec:deploy_app}

部署 Nginx 应用：

\begin{verbatim}
kubectl create deployment nginx --image=nginx
kubectl expose deployment nginx --port=80 --type=NodePort
\end{verbatim}

验证部署：
\begin{verbatim}
kubectl get deployments
kubectl get pods
kubectl get services
\end{verbatim}

\section{滚动更新}
\label{sec:rolling_update}

更新应用版本：
\begin{verbatim}
kubectl set image deployment/nginx nginx=nginx:1.24
kubectl rollout status deployment/nginx
\end{verbatim}

回滚更新：
\begin{verbatim}
kubectl rollout undo deployment/nginx
\end{verbatim}

\section{扩展应用}
\label{sec:scale_app}

手动扩展：
\begin{verbatim}
kubectl scale deployment/nginx --replicas=5
\end{verbatim}

自动扩缩容（HPA）：
\begin{verbatim}
kubectl autoscale deployment/nginx --min=2 --max=10 --cpu-percent=80
\end{verbatim}

\backmatter

\end{document}